
\documentclass[11pt]{article}

\usepackage[a4paper,margin=1in]{geometry}
\usepackage{amsmath,amssymb,amsthm,mathtools}
\usepackage[T1]{fontenc}
\usepackage{lmodern}
\usepackage{microtype}
\usepackage{hyperref}

\newtheorem{theorem}{Theorem}[section]
\newtheorem{proposition}[theorem]{Proposition}
\newtheorem{corollary}[theorem]{Corollary}
\newtheorem{conjecture}[theorem]{Conjecture}
\newtheorem{remark}[theorem]{Remark}
\newtheorem{definition}[theorem]{Definition}

\newcommand{\C}{\mathbb{C}}
\newcommand{\R}{\mathbb{R}}
\newcommand{\eps}{\varepsilon}
\newcommand{\sig}{\sigma}
\newcommand{\tridiag}{\operatorname{tridiag}}
\newcommand{\smin}{s_{\min}}
\newcommand{\conn}{\mathrm{conn}}
\newcommand{\scn}{\mathrm{sc}}
\newcommand{\dist}{\operatorname{dist}}
\newcommand{\diag}{\operatorname{diag}}
\newcommand{\I}{\mathrm{i}}

\title{A Note on the Simply Connectedness Threshold Problem\\
for the Clean OBC Hatano--Nelson Family}
\author{}
\date{}

\begin{document}
\maketitle

\begin{abstract}
We collect the core definitions, conjectures, and the strongest results currently established for the \emph{clean open-boundary Hatano--Nelson family}
\[
A_N=\tridiag(b,\mu,a),\qquad a>b>0,\ \mu\in\R.
\]
The main open problem is whether, once the strict $\eps$-pseudospectrum becomes connected, it is automatically simply connected. We record precise connectedness-threshold estimates, a strip confinement result for possible holes, an exact $N=2$ theorem, and two equivalent-looking stronger conjectures based on vertical slices and inertia monotonicity.
\end{abstract}

\section{Setup and strict pseudospectra}

We work with the real nonnormal tridiagonal Toeplitz matrices
\[
A_N=\tridiag(b,\mu,a)\in\C^{N\times N},
\qquad a>b>0,\quad \mu\in\R.
\]
Thus
\[
A_N=
\begin{pmatrix}
\mu & a \\
b & \mu & a \\
& \ddots & \ddots & \ddots \\
&& b & \mu & a \\
&&& b & \mu
\end{pmatrix}.
\]

Following Trefethen--Embree, we use the \emph{strict} $\eps$-pseudospectrum \cite{TE}:
\begin{definition}
For a matrix $A\in\C^{N\times N}$ and $\eps>0$,
\[
\sig_\eps(A):=\bigl\{z\in\C:\|(zI-A)^{-1}\|>\eps^{-1}\bigr\}
=\bigl\{z\in\C:\smin(zI-A)<\eps\bigr\}.
\]
\end{definition}

For matrices, $\sig_\eps(A)$ is an open subset of $\C$, each connected component of $\sig_\eps(A)$ contains at least one eigenvalue of $A$, and $\sig_\eps(A)$ has at most $N$ connected components \cite{TE}.

\begin{definition}
For the family $A_N$, define the \emph{connectedness threshold}
\[
\eps_N^{\conn}:=\inf\{\eps>0:\ \sig_\eps(A_N)\ \text{is connected}\},
\]
and the \emph{simply connectedness threshold}
\[
\eps_N^{\scn}:=\inf\{\eps>0:\ \sig_\eps(A_N)\ \text{is simply connected}\}.
\]
Equivalently, for fixed $\eps>0$ define
\[
N_{\conn}(\eps):=\min\{N\in\mathbb N:\ \sig_\eps(A_N)\ \text{is connected}\},
\]
\[
N_{\scn}(\eps):=\min\{N\in\mathbb N:\ \sig_\eps(A_N)\ \text{is simply connected}\},
\]
whenever these minima exist.
\end{definition}

\section{Basic explicit structure}

Set
\[
\rho:=\sqrt{b/a}\in(0,1),\qquad t:=\sqrt{ab}.
\]

\begin{proposition}[Similarity to a symmetric tridiagonal matrix]
\label{prop:similarity}
Let
\[
D_N:=\diag(1,\rho,\rho^2,\dots,\rho^{N-1}),
\qquad
S_N:=\tridiag(t,\mu,t).
\]
Then
\[
A_N=D_NS_ND_N^{-1}.
\]
Hence the eigenvalues of $A_N$ are
\[
\lambda_j=\mu+2t\cos\frac{j\pi}{N+1},
\qquad j=1,\dots,N.
\]
Moreover
\[
\kappa_2(D_N)=\rho^{-(N-1)}=(a/b)^{(N-1)/2}.
\]
\end{proposition}

\begin{remark}
The explicit eigenvalues and eigenvectors of tridiagonal Toeplitz matrices, together with sensitivity information and pseudospectral formulae, are classical; see \cite{NPR}. The ellipse-like geometry of tridiagonal Toeplitz pseudospectra goes back to Reichel--Trefethen \cite{RT}.
\end{remark}

\section{Main open problem}

\begin{conjecture}[Connected implies simply connected]
\label{conj:main}
For every $N\ge 2$ and every $\eps>0$,
\[
\sig_\eps(A_N)\ \text{connected}
\quad\Longrightarrow\quad
\sig_\eps(A_N)\ \text{simply connected}.
\]
Equivalently,
\[
\eps_N^{\scn}=\eps_N^{\conn}
\qquad\text{for all }N\ge 2,
\]
or, in inverse form,
\[
N_{\scn}(\eps)=N_{\conn}(\eps)
\qquad\text{for all }\eps>0.
\]
\end{conjecture}

\section{Best current connectedness-threshold estimates}

Define the exact nearest-neighbor spectral gap
\[
g_N:=\min_{1\le j\le N-1}|\lambda_j-\lambda_{j+1}|
=
4\sqrt{ab}\,
\sin\frac{\pi}{2(N+1)}\,
\sin\frac{3\pi}{2(N+1)},
\]
and
\[
C_N(\rho):=
\sup_{\theta\in(0,\pi)}
\frac{\bigl|\sin((N+1)\theta)\bigr|}
{\left(\sum_{j=1}^{N}\rho^{2(j-1)}\sin^2(j\theta)\right)^{1/2}}.
\]

\begin{theorem}[Connectedness threshold bounds]
\label{thm:conn-bounds}
The connectedness threshold satisfies
\[
\frac{g_N}{2}\,\rho^{N-1}
\;\le\;
\eps_N^{\conn}
\;\le\;
a\,\rho^N\,C_N(\rho).
\]
Equivalently:

\begin{itemize}
\item if
\[
\eps<\frac{g_N}{2}\rho^{N-1},
\]
then $\sig_\eps(A_N)$ is disconnected;

\item if
\[
\eps>a\rho^N C_N(\rho),
\]
then $\sig_\eps(A_N)$ is connected.
\end{itemize}

In particular,
\[
\eps_N^{\conn}\le \sqrt{ab}\,(N+1)\rho^{N-1},
\]
since $C_N(\rho)\le N+1$.
\end{theorem}

\begin{corollary}[Asymptotic connectedness threshold]
\label{cor:asymptotic-conn}
As $\eps\downarrow 0$,
\[
N_{\conn}(\eps)
=
\frac{2}{\log(a/b)}\log\frac1\eps
+
O\!\left(\log\log\frac1\eps\right).
\]
Equivalently,
\[
\eps_N^{\conn}
=
\rho^N \cdot \mathrm{poly}(N)
\]
with the sharp exponential rate $\rho^N=(b/a)^{N/2}$.
\end{corollary}

\section{Proved partial results toward simple connectedness}

\subsection{Outer-strip monotonicity}

For fixed $x\in\R$, define
\[
s_N(x,y):=\smin\bigl((x+\I y)I-A_N\bigr).
\]
Set
\[
\kappa_N:=(a-b)\cos\frac{\pi}{N+1}.
\]

\begin{proposition}[Monotonicity outside the central strip]
\label{prop:strip}
For every fixed $x\in\R$, the function $y\mapsto s_N(x,y)$ is even, and it is nondecreasing on $[\kappa_N,\infty)$.

Consequently, any bounded complementary component (hole) of $\sig_\eps(A_N)$ must lie in the horizontal strip
\[
|\Im z|<\kappa_N.
\]
\end{proposition}

\begin{remark}
This is a rigorous partial monotonicity statement. It does \emph{not} prove simple connectedness, but it shows that any counterexample must be confined to a narrow central strip.
\end{remark}

\subsection{The case \texorpdfstring{\(N=2\)}{N=2}}

\begin{proposition}[Exact \(N=2\) formula and simple connectedness]
\label{prop:N2}
For
\[
A_2=
\begin{pmatrix}
\mu & a\\
b & \mu
\end{pmatrix},
\]
if $d:=x-\mu$, then
\[
s_2(x,y)^2
=
d^2+y^2+\frac{a^2+b^2}{2}
-\frac12\sqrt{(a^2-b^2)^2+4(a+b)^2d^2+4(a-b)^2y^2}.
\]
In particular, for every fixed $x$,
\[
\frac{\partial}{\partial y}s_2(x,y)>0
\qquad (y>0).
\]
Hence every vertical section of $\sig_\eps(A_2)$ is an interval, and therefore
\[
\sig_\eps(A_2)\ \text{connected}
\quad\Longleftrightarrow\quad
\sig_\eps(A_2)\ \text{simply connected}.
\]
Equivalently,
\[
\eps_2^{\scn}=\eps_2^{\conn}.
\]
\end{proposition}

\section{Two useful reformulations of the open problem}

\subsection{The determinant route: a sufficient criterion}

Fix $x\in\R$ and write $u:=y^2\ge 0$. Define the Hermitian matrix
\[
H_{N,\eps}(x,u)
:=
\bigl((x+\I\sqrt u)I-A_N\bigr)^*
\bigl((x+\I\sqrt u)I-A_N\bigr)-\eps^2I,
\]
and the polynomial
\[
P_{N,\eps}(x,u):=\det H_{N,\eps}(x,u).
\]

\begin{proposition}[Root-counting criterion]
\label{prop:det-sufficient}
Fix $N$ and $\eps$. If for every $x\in\R$ the polynomial
\[
u\longmapsto P_{N,\eps}(x,u)
\]
has at most one positive root, then
\[
\sig_\eps(A_N)\ \text{connected}
\quad\Longrightarrow\quad
\sig_\eps(A_N)\ \text{simply connected}.
\]
\end{proposition}

\begin{remark}[Limitation of the determinant route]
Numerically, the criterion in Proposition~\ref{prop:det-sufficient} is \emph{sufficient but not necessary}. In scans near the connectedness threshold one can find examples where $P_{N,\eps}(x,\cdot)$ has two positive roots, even though the minimal branch
\[
m_{N,\eps}(x,y):=\lambda_{\min}\!\bigl(H_{N,\eps}(x,y^2)\bigr)
=
s_N(x,y)^2-\eps^2
\]
changes sign only once. Thus $\det H_{N,\eps}$ may overcount crossings coming from nonminimal singular-value branches.
\end{remark}

\subsection{The inertia-flow route: the strongest current conjectural reformulation}

Define
\[
\nu_{N,\eps}(x,y):=
n_-\!\bigl(H_{N,\eps}(x,y^2)\bigr),
\]
the number of negative eigenvalues of the Hermitian matrix \(H_{N,\eps}(x,y^2)\).

Then
\[
x+\I y\in \sig_\eps(A_N)
\quad\Longleftrightarrow\quad
\nu_{N,\eps}(x,y)\ge 1.
\]

\begin{conjecture}[Vertical inertia monotonicity]
\label{conj:inertia}
For every fixed \(N\), \(\eps\), and \(x\in\R\), the function
\[
y\longmapsto \nu_{N,\eps}(x,y)
\]
is nonincreasing on \([0,\infty)\).
\end{conjecture}

\begin{proposition}
Conjecture~\ref{conj:inertia} implies Conjecture~\ref{conj:main}.
\end{proposition}

\begin{remark}
Numerically, Conjecture~\ref{conj:inertia} appears significantly more robust than monotonicity of the minimal branch alone. In all scans performed so far (including near-threshold cases with \(a/b\) close to \(1\)), no upward jump of \(\nu_{N,\eps}(x,y)\) was observed as \(y\) increased.
\end{remark}

\section{Summary of the current state}

For the clean OBC Hatano--Nelson family \(A_N=\tridiag(b,\mu,a)\) with \(a>b>0\):

\begin{enumerate}
\item The strict connectedness threshold \(\eps_N^{\conn}\) is understood up to the correct exponential scale:
\[
\frac{g_N}{2}\rho^{N-1}\le \eps_N^{\conn}\le a\rho^N C_N(\rho),
\qquad
N_{\conn}(\eps)=\frac{2}{\log(a/b)}\log\frac1\eps+O(\log\log(1/\eps)).
\]

\item The main open question is whether connectedness already forces simple connectedness:
\[
\eps_N^{\scn}\stackrel{?}= \eps_N^{\conn}.
\]

\item This is known to hold for \(N=2\).

\item Any hole, if it exists, must lie in the strip
\[
|\Im z|<(a-b)\cos\frac{\pi}{N+1}.
\]

\item Two promising reformulations are:
\begin{itemize}
\item the polynomial root-counting problem for \(P_{N,\eps}(x,u)\);
\item the inertia-flow conjecture for \(\nu_{N,\eps}(x,y)\), which appears to be the sharper and more stable formulation.
\end{itemize}
\end{enumerate}

\begin{thebibliography}{99}

\bibitem{TE}
L.~N.~Trefethen and M.~Embree,
\emph{Spectra and Pseudospectra: The Behavior of Nonnormal Matrices and Operators},
Princeton University Press, Princeton, 2005.

\bibitem{RT}
L.~Reichel and L.~N.~Trefethen,
Eigenvalues and pseudo-eigenvalues of Toeplitz matrices,
\emph{Linear Algebra Appl.} \textbf{162--164} (1992), 153--185.

\bibitem{NPR}
S.~Noschese, L.~Pasquini, and L.~Reichel,
Tridiagonal Toeplitz matrices: properties and novel applications,
\emph{Numer. Linear Algebra Appl.} \textbf{20} (2013), 302--326.

\end{thebibliography}

\end{document}
